% ===================================================================
%  اللوحة رقم ١٠ — البحر. — الميناء.
%  Traduction 2026 d'apres texto/io, controlee sur texto/fr, texto/en
%  et texto/es. Consigne : voir texto/ar/10-lawha-01.tex.
%
%  Le francais decale sa numerotation d'un cran a partir de l'alinea 8.
%  C'est une coquille de l'imprimeur francais, conservee dans la
%  transcription ; l'arabe suit l'ido, qui compte juste.
% ===================================================================

%%K t10-apar-1 apar
\VUsaut{4.0mm}
\VUtitre{15.0pt}{60}{اللوحة رقم ١٠}
\VUsaut{3.0mm}
\VUpk{11.4pt}{40}{البحر. --- الميناء.}
\VUsaut{3.0mm}

%%K t10-01-1 p
1. --- في الصيف، بينما يطلب بعضهم الهدوء والهواء البارد في الجبال، يؤثر
آخرون الذهاب إلى الساحل. وذلك ما فعلناه العام الماضي، فنحن نحب
\VUgras{المحيط} \textsuperscript{(1)} حبًا شديدًا.

%%K t10-02-1 p
2. --- أما أنا فأجد لذة كبيرة في التحديق في تلك الصفحة الهائلة من الماء
الممتدة إلى \VUgras{الأفق} \textsuperscript{(2)}، وفي مشاهدة \VUgras{البواخر}
\textsuperscript{(3)} الضخمة تنطلق في \VUgras{عبور} طويل وعلى متنها المسافرون
إلى أمريكا.

%%K t10-03-1 p
3. --- ولذلك يختار أبي، وهو يعرف ما أحب، للعطلة دائمًا شاطئًا هادئًا جدًا
قرب أحد \VUgras{موانينا الحربية} الكبيرة.

%%K t10-03-2 p
وفي إقامتنا الأخيرة جاء \VUgras{الأسطول} فألقى مراسيه خلف \VUgras{حاجز
الأمواج} \textsuperscript{(4)} الضخم الذي يغلق \VUgras{الحوض الحربي}
\textsuperscript{(5)} ويحميه، فتسنّى لي أن أتأمل على مهل مختلف \VUgras{السفن
الحربية}: من \VUgras{الغواصة} \textsuperscript{(10)} الصغيرة التي تغوص وتختفي
تحت الماء، إلى \VUgras{البارجة} \textsuperscript{(9)} الهائلة التي لم يكن
يخيفها حتى الآن إلا هجمات \VUgras{زورق الطوربيد} \textsuperscript{(8)}.

%%K t10-tit-2 sub
\VUsaut{3.0mm}
\VUpk{10.2pt}{40}{السفن الصغيرة.}
\VUsaut{3.0mm}

%%K t10-04-1 p
4. --- وقد كان هذا الأخير يستطيع من قبلُ أن يجد، في \VUgras{مهارة}،
الموضع الضعيف في الدرع المعدني الثخين فيغرس فيه \VUgras{الطوربيد} الخطر
الذي يُغرِق انفجارُه السفينةَ؛ لكنه كان مع ذلك أقل مدعاةً للخوف من
الغواصة، إذ تستطيع المدمّرات مطاردته بسهولة.

%%K t10-05-1 p
5. --- ورأيت أيضًا \VUgras{الطرّادات} \textsuperscript{(6)} المدرّعة السريعة،
وهي تكاد تكون منيعة كبارجة حقيقية، وعدة \VUgras{سفن نقل جند}
\textsuperscript{(12)}، وثلاثة أو أربعة \VUgras{زوارق مدفعية} \textsuperscript{(7)}،
وأخيرًا \VUgras{فرقاطة} \textsuperscript{(11)}. \VUgras{ومنظر ذلك الأسطول وهو
يرفع مراسيه أطرف ما في الأمر كله.}

%%K t10-06-1 p
6. --- وما أكبر الفرق في البنية بين تلك الكتل الفولاذية الضخمة وبين
\VUgras{السفن ذوات الصواري الثلاثة} الأنيقة أو \VUgras{المراكب الشراعية}
\textsuperscript{(13)} الرشيقة التي لا تزال تُستعمل في الأسطول التجاري، والتي
كنت أراها تدخل الميناء ناشرةً كل قلوعها وأنا أتمشى على \VUgras{الرصيف}
\textsuperscript{(14)}. صحيح أنها كانت تفقد كثيرًا من ميزتها إذا كانت
\VUgras{الريح} ضدها. فقد كان عليها أن تطوي كل قلوعها لتدخل الحوض، وكان
لا بد من \VUgras{قاطرة} \textsuperscript{(16)} تأتي فتجرّها إلى الرصيف
كـ\VUgras{الصنادل} \textsuperscript{(17)} البسيطة.

%%K t10-tit-3 sub
\VUsaut{3.0mm}
\VUpk{10.2pt}{40}{الميناء التجاري.}
\VUsaut{3.0mm}

%%K t10-07-1 p
7. --- والذي يجعل الميناء الذي أتحدث عنه بهذه الطرافة أنه لا يضم حوضًا
حربيًا صُنع صناعةً بأشغال جبّارة فحسب، بل يضم أيضًا \VUgras{ميناءً تجاريًا}
رائعًا عند \VUgras{مصبّ} نهر كبير.

%%K t10-08-1 p
8. --- ولسان الأرض الفاصل بين الحوضين محصَّن، تدافع عنه سلسلة من
\VUgras{البطاريات} و\VUgras{الحصون} الممتدة في البحر. ولا بد من إذن خاص
للتمشي على \VUgras{الأسوار} \textsuperscript{(15)}. وقد نلته بسهولة عن طريق
عمي، وهو مهندس في \VUgras{أحواض بناء السفن} \textsuperscript{(18)}، فذهبت أرى
السفن تُبنى تحت سقائف واسعة.

%%K t10-09-1 p
9. --- بل رأيت إنزال بارجة إلى الماء، وأذكر تلك الرعشة التي سرت في الحشد
كله حين بدأت الكتلة الهائلة، وقد تُركت لنفسها، تنزلق على مهدها فطفت على
ماء النهر.

%%K t10-10-1 p
10. --- وللوصول إلى أحواض البناء تعبر \VUgras{ساحة العرض} \textsuperscript{(19)}
حيث تأتي حامية المدينة للتدريب كل يوم، وتمرّ على \VUgras{جسر مشاة}
\textsuperscript{(20)} حديدي يصل ساحة العرض بـ\VUgras{الترسانة} \textsuperscript{(21)}.

%%K t10-tit-4 sub
\VUsaut{3.0mm}
\VUpk{10.2pt}{40}{الترسانة والأحواض.}
\VUsaut{3.0mm}

%%K t10-11-1 p
11. --- وزرت مع عمي تلك المنشأة الكبيرة الملأى بـ\VUgras{المدافع}
\textsuperscript{(22)} و\VUgras{القذائف} \textsuperscript{(23)} و\VUgras{القنابل}.
ورأيت أيضًا \VUgras{مصنع الحديد} \textsuperscript{(24)} حيث تُصنع
\VUgras{مراجل} \textsuperscript{(25)} البواخر.

%%K t10-12-1 p
12. --- وعلى مقربة \VUgras{الأحواض} \textsuperscript{(26)} --- الأحواض الجافة ---
و\VUgras{الأحواض العائمة} \textsuperscript{(27)} اللافتة جدًا، حيث تسنّى لي أن
أرى عن قرب، على سفينة تحت الإصلاح، \VUgras{مروحة} \textsuperscript{(28)} السفينة
و\VUgras{عارضتها} \textsuperscript{(29)} و\VUgras{دفّتها} \textsuperscript{(30)}.

%%K t10-13-1 p
13. --- والميناء التجاري جدير بالدرس كالحربي تمامًا، وقد قضيت ساعات كثيرة
أتسكّع على الأرصفة حيث ترسو \VUgras{البواخر ذوات العجلات} \textsuperscript{(31)}
الصاعدة في النهر، وأرقب \VUgras{الرافعة ثلاثية القوائم} \textsuperscript{(32)}
وهي تعمل، ترفع أحمالًا هائلة كأنها ريش.

%%K t10-tit-5 sub
\VUsaut{4.0mm}
\VUpk{10.2pt}{40}{المرشد البحري.}
\VUsaut{3.0mm}

%%K t10-14-1 p
14. --- وأعجبت بـ\VUgras{البواخر العابرة للأطلسي} \textsuperscript{(34)} تغادر
المرسى و\VUgras{أعلامها} \textsuperscript{(35)} خافقة، يقودها \VUgras{مرشد}
\textsuperscript{(36)} ماهر يعرف عجيب المعرفة كيف يتبع \VUgras{المجرى} المعلَّم
بـ\VUgras{العوّامات} \textsuperscript{(40)} و\VUgras{المنارات}، متجنبًا
\VUgras{صنادل النقل} \textsuperscript{(41)} التي يعسر توجيهها بقلعها المربّع
الوحيد. وكنت أتبيّن عند \VUgras{المقدّم} \textsuperscript{(37)} \VUgras{مقصورات}
\textsuperscript{(39)} الدرجة الثانية، وعند \VUgras{المؤخّر} \textsuperscript{(38)}
صالات ركاب الدرجة الأولى.

%%K t10-15-1 p
15. --- وأحيانًا كنت أرقب شحن \VUgras{سفينة بضائع} \textsuperscript{(42)} تُكدَّس
فيها البضائع الآتية من \VUgras{المخزن} \textsuperscript{(43)} و\VUgras{محطة
الميناء} \textsuperscript{(44)}، تجلبها قطارات \VUgras{خط الرصيف}
\textsuperscript{(45)} الكثيرة.

%%K t10-tit-6 sub
\VUsaut{3.0mm}
\VUpk{10.2pt}{40}{التجديف للمتعة.}
\VUsaut{3.0mm}

%%K t10-16-1 p
16. --- وكثيرًا ما ركبت الزوارق في \VUgras{الحوض المغلق} \textsuperscript{(53)}
حيث تحتمي \VUgras{القوارب} \textsuperscript{(46)} الصغيرة، وكان مسك
\VUgras{المجداف} \textsuperscript{(47)} فرحة لي. كنت أصعد إلى \VUgras{سطح}
\textsuperscript{(48)} \VUgras{قارب شراعي} \textsuperscript{(49)}، وأتسلق
عبر \VUgras{الحبال} \textsuperscript{(52)} \VUgras{الصاري}
\textsuperscript{(51)} إلى \VUgras{العوارض} \textsuperscript{(50)}, ثم أتابع نزهتي في \VUgras{زورق}
\textsuperscript{(54)} بين \VUgras{قوارب الصيد} \textsuperscript{(55)} الثقيلة حيث
تُنشر \VUgras{الشباك} \textsuperscript{(56)} لتجف.

%%K t10-17-1 p
17. --- وعلى \VUgras{الرصيف} \textsuperscript{(58)} كانت تمر أمامي
\VUgras{البضائع} \textsuperscript{(59)} على اختلافها، معبّأة في \VUgras{صناديق}
\textsuperscript{(60)} أو في \VUgras{بالات} \textsuperscript{(61)}. وكانت تؤلف
\VUgras{حمولة} \textsuperscript{(63)} الصنادل، فترفعها \VUgras{رافعات}
\textsuperscript{(62)} قوية وتضعها على \VUgras{شاحنات} \textsuperscript{(64)}، بينما
يصرّح عنها \VUgras{الناقلون} ويؤدّون الرسم في \VUgras{الجمرك}
\textsuperscript{(65)}.

%%K t10-18-1 p
18. --- وأخيرًا، متى بلغت \VUgras{الجسر المتحرك} \textsuperscript{(66)} أو
\VUgras{الهويس} \textsuperscript{(69)}، مارًّا بـ\VUgras{الكرّاءة}
\textsuperscript{(68)} التي تنظف \VUgras{القناة} \textsuperscript{(67)}، كنت أنزل
على الرصيف قرب \VUgras{صاري الإشارات} \textsuperscript{(70)}.

%%K t10-19-1 p
19. --- وعلى الجانب الآخر من ذلك الرصيف ترسو \VUgras{الزوارق البخارية}
\textsuperscript{(71)} التي تحمل ضباط الأسطول إلى البر. ومنظر بديع أن ترى
البحّارة يرفعون مجاديفهم على إيقاع واحد بينما يقترب الزورق، تحمله
\VUgras{الموجة} \textsuperscript{(72)}، من درج المرسى في يسر. وهنا يمكن أن
يُرقَب \VUgras{المدّ} وحركة \VUgras{الجزر} و\VUgras{المدّ} على
\VUgras{الشاطئ} \textsuperscript{(73)} خير رقبة.

%%K t10-tit-7 sub
\VUsaut{3.0mm}
\VUpk{10.2pt}{40}{نادي الضباط.}
\VUsaut{3.0mm}

%%K t10-20-1 p
20. --- وللعودة إلى البيت كنت أركب \VUgras{الترام الكهربائي}
\textsuperscript{(74)}، و\VUgras{موقفه} \textsuperscript{(75)} عند \VUgras{العمود}
القائم أمام نادي الضباط.

%%K t10-20-2 p
ومن \VUgras{شرفة} \textsuperscript{(76)} النادي، المزينة بـ\VUgras{مراسٍ}
\textsuperscript{(77)} ومدافع وقذائف، كثيرًا ما رأيت جمعًا بهيًا من ضباط
البحرية: \VUgras{ملازمين} \textsuperscript{(78)} بسيوفهم، و\VUgras{نقباء مدفعية}
\textsuperscript{(79)} و\VUgras{مشاة} \textsuperscript{(80)} بـ\VUgras{سيوفهم}.
وخلفهم كان يقف \VUgras{بحّار} \textsuperscript{(81)} يناولهم \VUgras{منظارًا}
متى أرادوا أن يتبيّنوا ما يجري في البعد.

%%K t10-21-1 p
21. --- ورأيت أيضًا \VUgras{ملازمين ثوانيَ} \textsuperscript{(82)} شبابًا يتلقّون
أوامرهم من \VUgras{قائدهم} \textsuperscript{(83)} أو من \VUgras{الأميرال}
\textsuperscript{(84)}، بينما ينتظر \VUgras{رئيس عرفاء} \textsuperscript{(85)}
ليحملها إلى السفينة.

%%K t10-22-1 p
22. --- ومن تلك الشرفة المنظر رائع: فهي تشرف على الشاطئ كله
بـ\VUgras{خيامه} \textsuperscript{(86)} و\VUgras{أكشاك الاستحمام} \textsuperscript{(87)}،
وفي ساعة الاستحمام يُرى \VUgras{المستحمّون} \textsuperscript{(88)} يمرحون في
بهجة أمام \VUgras{الكازينو} \textsuperscript{(89)}.

%%K t10-23-1 p
23. --- وفي آخر الشاطئ يؤلف الساحل \VUgras{شبه جزيرة} \textsuperscript{(91)}
صخرية صغيرة، تتحطم عليها \VUgras{الأمواج المتكسّرة} \textsuperscript{(92)}،
وقد بُنيت عليها \VUgras{منارة} \textsuperscript{(93)} تدل البحّارة على طريقهم
ليلًا. وشبه الجزيرة هذه لا تتصل بالبر إلا بـ\VUgras{برزخ} \textsuperscript{(90)}
شديد الضيق.

%%K t10-tit-8 sub
\VUsaut{3.0mm}
\VUpk{10.2pt}{40}{منظر عام للساحل.}
\VUsaut{3.0mm}

%%K t10-24-1 p
24. --- ويمتد \VUgras{الجرف} \textsuperscript{(94)} وقد نحته البحر فحفر فيه
سلسلة عجيبة من \VUgras{الخلجان} \textsuperscript{(95)} و\VUgras{الرؤوس}، فجعله
في غاية الجمال.

%%K t10-24-2 p
وعلى \VUgras{الهضبة} \textsuperscript{(96)} المشرفة على \VUgras{المضيق}
\textsuperscript{(98)} \VUgras{حصن} \textsuperscript{(97)} بالغ الأهمية وبضع فيلات.

%%K t10-25-1 p
25. --- وذلك المضيق يفصل عن البر \VUgras{جزيرة} \textsuperscript{(99)} شديدة
الخطر، تحفّ بها \VUgras{الصخور} \textsuperscript{(100)} و\VUgras{الشعاب}، وفيها
تتحطم أحيانًا القوارب بل السفن الكبيرة. ومهما كانت النجدة سريعة ومهما
عجّلت \VUgras{قوارب النجاة}، فإن هذه \VUgras{حوادث الغرق} مروّعة، وقد
هلكت في \VUgras{عواصف} الاعتدال عدة سفن بمن فيها وما فيها.

%%K t10-25-2 p
ورغم هذا الجوار، فإن \VUgras{الميناء} كثير الارتياد من البحّارة.
\VUsaut{6.0mm}
\VUfilet{22mm}
